\documentclass[10pt,a4paper,landscape,oneside]{amsart}
\usepackage{amsmath, amsthm, amssymb, amsfonts}
\usepackage{graphicx}
\usepackage[T1]{fontenc}
\usepackage{booktabs}
\usepackage{color}
\usepackage{fancyhdr}
\usepackage{float}
\usepackage{fullpage}
\usepackage[scaled]{beramono}
\usepackage{enumitem}
\usepackage{multicol}
\usepackage{fix-cm}

\usepackage{minted}
\newcommand{\code}[1]{\inputminted[fontsize=\large,tabsize=2,baselinestretch=1]{java}{#1}}

\pagestyle{fancy}
\lhead{The Floor is Made of Java, University of Copenhagen}
\rhead{\thepage}
\chead{}
\setlength{\headheight}{15.2pt}

\begin{document}
\pagestyle{empty}
\pagenumbering{gobble}

\begin{center}
	\strut
	\vspace{1cm}\\
	\includegraphics[width=3cm]{ku.eps}
	\vspace{0.2cm}\\
	{\Huge University of Copenhagen\\}
	{\fontsize{40}{40}\selectfont The Floor is Made of Java\\}
	\vspace{0.5cm}
	{\Large Alice Ryhl, Carl Dybdahl, Jonas Klausen\\}
	\vfill
	{\huge NWERC 2019}\\
	\vspace{1cm}
	{\Large \today}\\
	\vspace{1cm}
\end{center}
\clearpage

\pagestyle{fancy}
\pagenumbering{arabic}

\begin{multicols*}{2}
\section{Number theory}
\subsection{Euclid}
Computes \(x\), \(y\) and \(d\) such that \(xa + yb = d = \operatorname{gcd}(a,b)\). \\
Note that you can find \texttt{lcm} using \(\operatorname{lcm}(a,b) = \frac{a \cdot
b}{\operatorname{gcd}(a,b)}\).
\code{number-theory/euclid.java}
\subsection{Modular inverse}
Assumes that \(a\) and \(mod\) are relatively prime.
\code{number-theory/mod-inverse.java}
\subsection{Modular power}
Computes \texttt{a\textsuperscript e \% mod}.
\code{number-theory/mod-pow.java}
\subsection{Chinese remainder theorem}
The \texttt{crt2} function finds \(x\) with \(x \equiv a \mod m\) and \(x \equiv b \mod
n\), while the \texttt{crt} function combines it to solve it for any number of equations.
Returns \texttt{true} if it found a solution.
\code{number-theory/crt.java}
\subsection{Prime sieve}
Compute all primes up to some number. Runs in \(O(n \log n \log \log n)\).
\code{number-theory/sieve.java}
\subsection{Binomial coefficients}
Compute the binomial coefficients.
\code{number-theory/binom.java}
\subsection{Catalan numbers}
Compute the \(n\) first Catalan numbers. These can be used for various things:
\begin{enumerate}
  \item Number of expressions containing \(n\) pairs of parentheses which are correctly
    matched.
  \item Number of full binary trees with \(n+1\) leaves.
  \item Number of lattice paths in \(n \times n\) square that stay below the diagonal.
  \item A convex polygon with \(n+2\) sides can be cut into \(n\) triangles in \(C_n\)
    ways.
  \item Number of permutations with no three-term increasing subsequence.
\end{enumerate}
This uses the recursive formula \(C_0 = 1\) and \(C_{n+1} = \frac{2(2n+1)}{n+2} C_n\).
Could also use \(C_n = \frac{1}{n+1}\binom{2n}{n} = \binom{2n}{n} - \binom{2n}{n+1}\).
\code{number-theory/catalan.java}

\section{Graph algorithms}
\subsection{Shortest path, negative loops}
A variant of Bellman Ford that computes shortest path. Marks unreachable nodes with
\verb|MAX_VALUE| and nodes reachable from a negative loop with \verb|MIN_VALUE|.
If there are no negative weight
\code{graph/shortest-path.java}
\subsection{Shortest path, no negative loops}
A simpler version of Bellman Ford that never terminates when given a negative loop.
\code{graph/shortest-path-noloop.java}
\subsection{Floyd Warshall}
Shortest path between all pairs of verticies. Negative weights are allowed. Runs in
\(O(n^3)\).
\code{graph/floyd-warshall.java}

\section{Geometry}
\subsection{Convex hull}
Computes the convex hull in counter clockwise order in \(O(n \log n)\) time.
Use \verb|>=| instead of \verb|<=| to get clockwise order.
\code{geometry/convex-hull.java}
\subsection{Closest pair}
Find the closest pair of points in \(O(n (\log n)^2)\) time.
\code{geometry/closest-pair.java}
\subsection{Polygon area}
Finds the area of a polygon using the shoelace formula. Returns twice the signed area,
which is negative if points are in clockwise order.
\code{geometry/polygon-area.java}

\section{Linear algebra}
\subsection{Solve linear equations}
Find any solution of \(Ax = b\) in \(O(n^2m)\) time. Returns rank of solution, or \(-1\)
if no solution is found. Data in \texttt{A} and \texttt{b} is lost.
\code{linear-algebra/linear-system.java}
\subsection{Partial linear equations}
To get all uniquely determined values back, make the following changes. Non-unique values
are left as NaN.
\code{linear-algebra/partial-linear-system.java}
\end{multicols*}

\end{document}
% vim:spelllang=en
